% !TEX root = ../main.tex
\chapter[Búsqueda de hiperparámetros del imputador]{Búsqueda de Hiperparámetros del Imputador Supervisado}
\label{anx:tuning_imputador}

La estrategia de imputación supervisada descrita en la Sección~\ref{sec:demanda_censurada} se
apoya en un regresor LightGBM auxiliar. Este anexo documenta la búsqueda de hiperparámetros de
ese modelo, que constituye un procedimiento independiente del ajuste del modelo de
\textit{forecasting} y que se ejecuta como un modo de ejecución propio del sistema
(\texttt{run\_mode = tune\_imputation}). Se recoge aquí, y no en el cuerpo de la memoria, porque
su resultado no modifica la configuración del sistema final: lo que aporta es una lección
metodológica sobre la validez de una búsqueda de hiperparámetros, y el material que la sostiene.

\section{Protocolo}

La dificultad de fondo es que la demanda latente de un día con rotura real es un contrafactual
no observable, de modo que no existe verdad-terreno contra la que medir. La búsqueda opera por
tanto sobre el mismo protocolo de censura sintética de la
Sección~\ref{sec:imputacion_censura_sintetica}: se censura artificialmente una fracción de los
días limpios, se reconstruyen y se compara la estimación con la demanda conocida.

El espacio de búsqueda comprende \textbf{trece hiperparámetros}, con cotas acotadas en la
Tabla~\ref{tab:tuning_espacio}, recorridos por 300 \textit{trials} del muestreador gaussiano de
Optuna. Dos decisiones de diseño gobiernan el procedimiento:

\begin{itemize}
    \item \textbf{Promediado sobre múltiples extracciones.} En una sola extracción de censura
    sintética, la dispersión del error entre \textit{trials} (entre el 12\% y el 40\% según la
    medición) es un orden de magnitud mayor que las diferencias reales entre conjuntos de
    hiperparámetros (en torno al 1\%). Una búsqueda de extracción única selecciona por tanto el
    \textit{trial} afortunado, y su propio valor objetivo no es evidencia de mejora alguna: una
    ejecución que reportó un $-1{,}4\%$ en la muestra de selección midió $-0{,}32\%$, con
    intervalo de confianza cruzando el cero, sobre extracciones nuevas. El objetivo promedia
    en consecuencia sobre 15 extracciones.
    \item \textbf{Series disjuntas entre selección y validación.} Las semillas de censura
    disjuntas no bastan: cada extracción censura el 30\% del mismo conjunto de días limpios, de
    modo que tras $n$ extracciones el solapamiento es de $1 - 0{,}7^n$ y los dos conjuntos
    convergen sobre las mismas filas (se midió un 82{,}4\% de solapamiento con 5 y 10
    extracciones). La separación se hace por tanto sobre series: 35 series para seleccionar y
    15 reservadas para validar, con 25 extracciones sobre estas últimas.
\end{itemize}

El ganador solo se persiste si el intervalo de confianza al 95\% de la mejora media sobre las
extracciones de validación queda \textbf{íntegramente por debajo de cero}, y además solo si bate
tanto a los hiperparámetros por defecto como al campeón que ya estuviera en disco.

\begin{table}[htbp]
    \centering
    \small
    \begin{tabular}{llll}
        \toprule
        \textbf{Hiperparámetro} & \textbf{Cotas} & \textbf{Hiperparámetro} & \textbf{Cotas} \\
        \midrule
        \texttt{n\_estimators} & 50 -- 3.000 & \texttt{subsample} & 0,4 -- 1,0 \\
        \texttt{max\_depth} & 2 -- 12 & \texttt{learning\_rate} & 0,005 -- 0,3 (log) \\
        \texttt{num\_leaves} & 2 -- 1.024 & \texttt{colsample\_bytree} & 0,3 -- 1,0 \\
        \texttt{min\_child\_samples} & 2 -- 100 & \texttt{reg\_alpha} & $10^{-8}$ -- 100 (log) \\
        \texttt{min\_data\_per\_group} & 1 -- 100 & \texttt{reg\_lambda} & $10^{-8}$ -- 100 (log) \\
        \texttt{max\_bin} & 8 -- 255 & \texttt{cat\_smooth} & 0 -- 50 \\
        \texttt{subsample\_freq} & 0 -- 1 & & \\
        \bottomrule
    \end{tabular}
    \caption{Espacio de búsqueda de la sintonización del imputador supervisado. La cota superior
    de \texttt{n\_estimators} responde a una restricción de cómputo y está justificada por
    medición: a 500 series un ajuste de 8.000 árboles cuesta unos 34 segundos, lo que situaría
    una búsqueda de 300 \textit{trials} en torno a las 43 horas. Lo que se sacrifica al recortar
    a 3.000 son 0{,}0037 de MAE, cifra inferior a los 0{,}0059 de anchura del intervalo de
    validación que decide la persistencia, de modo que el recorte no puede alterar el veredicto.}
    \label{tab:tuning_espacio}
\end{table}

\section{Resultado: la ganancia se invierte con la escala}

Sobre el protocolo descrito, la búsqueda encuentra una configuración que mejora el error de
reconstrucción un \textbf{5{,}33\%} frente a los hiperparámetros por defecto, con un intervalo
de confianza al 95\% de $[-0{,}0199,\ -0{,}0120]$, enteramente por debajo de cero.

Ese resultado, sin embargo, no es transferible. Al repetir la medición variando únicamente el
número de días limpios con los que se ajusta el imputador, y manteniendo todo lo demás
constante, la ganancia se reduce de forma monótona y termina cambiando de signo
(Tabla~\ref{tab:tuning_reversion}).

\begin{table}[htbp]
    \centering
    \small
    \begin{tabular}{rlr}
        \toprule
        \textbf{Filas limpias de ajuste} & \textbf{Periodo de evaluación} & \textbf{Ganancia} \\
        \midrule
        467 (15 series reservadas) & series disjuntas & $-5{,}33\%$ \\
        595 (últimos 30 días de 50 series) & tiempo disjunto & $-4{,}34\%$ \\
        1.930 (50 series completas) & tiempo disjunto & $-1{,}72\%$ \\
        14.243 (500 series) & tiempo disjunto & $\mathbf{+12{,}44\%}$ \\
        18.310 (500 series más holdout) & semana de holdout & $\mathbf{+13{,}20\%}$ \\
        \bottomrule
    \end{tabular}
    \caption{Efecto del tamaño del conjunto de ajuste del imputador sobre la ganancia de la
    sintonización. Signo negativo indica mejora. Las dos últimas filas coinciden sobre periodos
    de evaluación distintos, lo que descarta el calendario como explicación y señala a la escala.}
    \label{tab:tuning_reversion}
\end{table}

La lectura es que \textbf{el tamaño del conjunto con el que se ajusta el imputador no es un
detalle del montaje experimental, sino la variable que decide la respuesta}. Los mismos
hiperparámetros que mejoran un 5\% con 467 filas de ajuste resultan un 12\% \emph{peores} que no
sintonizar a la escala de 500 series, que es precisamente la escala a la que el sistema ejecuta
la comparación de estrategias del Capítulo~\ref{chap:resultados}. La causa es reconocible: un
conjunto de ajuste pequeño premia configuraciones fuertemente regularizadas que dejan de ser
ventajosas cuando hay datos suficientes.

De ahí se derivan dos consecuencias prácticas. La primera es que un fichero de hiperparámetros
sintonizados solo es válido en el entorno del tamaño de ajuste con el que se obtuvo, y esa
condición no puede expresarse en el propio fichero, por lo que se registra explícitamente en los
metadatos de la búsqueda. La segunda, y es la que justifica la decisión del sistema final, es
que \textbf{las cifras de reconstrucción del Capítulo~\ref{chap:resultados} se obtienen con los
hiperparámetros por defecto}, no con los sintonizados: a la escala a la que se publican, los
defectos son la mejor opción disponible.

\section{Validación del buscador}

Un resultado como el anterior invita a una objeción legítima: que la reversión sea un artefacto
del muestreador empleado y no una propiedad del problema. Para acotarla se repitió la búsqueda
completa con los dos muestreadores disponibles, el basado en procesos gaussianos y el
\textit{Tree-structured Parzen Estimator}, sobre tres semillas comunes y 300 \textit{trials}
cada uno, en un total de unas 14 horas de cómputo (Tabla~\ref{tab:tuning_samplers}).

\begin{table}[htbp]
    \centering
    \small
    \begin{tabular}{rrrrrr}
        \toprule
        \textbf{Semilla} & \textbf{GP} & \textbf{TPE} & \textbf{Diferencia} & \textbf{GP (min)} & \textbf{TPE (min)} \\
        \midrule
        7 & $-6{,}32\%$ & $-6{,}69\%$ & $+0{,}37$ pp & 116 & 198 \\
        42 & $-7{,}62\%$ & $-4{,}56\%$ & $-3{,}06$ pp & 86 & 121 \\
        1.234 & $-6{,}29\%$ & $-7{,}48\%$ & $+1{,}19$ pp & 127 & 177 \\
        \midrule
        \textbf{Media} & $\mathbf{-6{,}74\%}$ & $\mathbf{-6{,}24\%}$ & $\mathbf{-0{,}50}$ \textbf{pp} & \textbf{110} & \textbf{165} \\
        \bottomrule
    \end{tabular}
    \caption{Comparación pareada de muestreadores sobre semillas comunes. La diferencia media
    entre ambos (0{,}50 puntos porcentuales) es cuatro veces menor que su desviación entre
    semillas (2{,}26 puntos), y el signo se invierte según la semilla.}
    \label{tab:tuning_samplers}
\end{table}

La conclusión es un \textbf{resultado nulo}: con tres semillas, el ruido entre extracciones
domina por completo la elección de muestreador, de modo que ambos conducen al mismo sitio y la
reversión de la sección anterior no es atribuible a ninguno de los dos. Lo que sí resulta
concluyente, y es la razón por la que el sistema emplea el muestreador gaussiano, es que este
resulta un 34\% más rápido y gana en las tres semillas sin excepción.

Dos observaciones adicionales del mismo experimento. Las seis búsquedas baten a los
hiperparámetros por defecto con el intervalo de confianza íntegramente por debajo de cero, entre
el $-4{,}56\%$ y el $-7{,}62\%$, lo que confirma que a esa escala de ajuste la ganancia es real
con independencia del muestreador. Y las configuraciones ganadoras \emph{no se parecen entre sí}:
\texttt{num\_leaves} recorre de 40 a 219, \texttt{n\_estimators} de 687 a 3.000 y
\texttt{learning\_rate} de 0{,}015 a 0{,}052, produciendo mejoras casi idénticas. La superficie
de respuesta es plana en la región buena, que es la misma observación con la que se justificó
promediar sobre múltiples extracciones.
